\hmsubsection{Hardware Integration}{UMA}{}
This section describes how the hardware components introduced in 
Section~7.2 are physically integrated into a working system. The 
description covers the data and power connections between the host 
processor, the motor controller, the camera, and the actuators, the 
configuration of the Dynamixel servos, and the runtime measures used 
to protect the elbow motor against thermal overload during 
prolonged operation.

\hmsubsubsection{Physical Connections}{UMA}{}
The Raspberry Pi~5 acts as the central host of the system and is 
connected to two peripheral devices through USB. The OAK-D~S2 
camera connects directly to one of the Pi's USB~3.0 ports through 
a USB-A to USB-C cable, providing both data transfer and the 
camera's required 5~V power. The OpenRB-150 motor controller connects to a second USB~3.0 port 
through a USB-C to USB-C cable, which carries serial data at 
115\,200~bps. The 12~V power supply for the 
motors is delivered to the OpenRB-150 through its dedicated motor 
power input, separately from the USB connection to the host. The 
Raspberry Pi itself is powered through a third USB-C cable from a 
regulated 5~V supply. All three power and data paths share a 
common ground reference.

The five Dynamixel servos are connected to the OpenRB-150 through 
the two-branch star topology described in Section~7.2.2 
(Figure~\ref{fig:power_topology}). Each branch leaves the controller 
through a separate 3-pin TTL cable, and the servos within each 
branch are linked through short JST cables in a daisy chain. The 
TTL bus carries both the 12~V motor power and the differential 
data signal on the same connector, in accordance with the 
Dynamixel Protocol 2.0 specification \cite{robotis_dynamixel}.

\hmsubsubsection{Motor Configuration}{UMA}{}
Before the servos can operate as a coordinated arm, each motor 
must be configured individually with a unique identifier on the 
TTL bus, a fixed baud rate, and a defined zero position. This 
configuration was performed using the ROBOTIS Dynamixel Wizard 
2.0 software with each motor connected to the OpenRB-150 in 
isolation. The five motors were assigned identifiers 1 through 5, 
corresponding to the base, shoulder, elbow, wrist, and gripper 
joints respectively. All five motors were set to a baud rate of 115\,200~bps, matching 
the baud rate used on the USB serial connection between the host 
processor and the OpenRB-150. Using a single rate on both sides 
of the controller simplifies the firmware and ensures that the 
limiting factor in the command-response cycle is the motor's 
mechanical response time rather than any rate mismatch between 
the two buses. While the Dynamixel servos support baud rates up 
to 4.5~Mbps, 115\,200~bps provides sufficient throughput for the 
movement command rate produced by the host, since the system 
issues at most a few commands per second during a typical 
pick-and-place cycle. Each motor 
was rotated to its mechanical centre and the corresponding encoder 
value was registered as the zero position, after which the joint 
limits described in Section~8.7 were applied to prevent 
the motors from rotating into mechanical hard stops during 
operation.

Two of the motors used in the project initially exhibited 
hardware error flags that prevented them from accepting commands. 
The flags were investigated by reading the hardware error status 
register from each motor over the TTL bus and were found to be 
latched overload errors from previous use. A factory reset was 
performed on the affected motors through the Dynamixel Wizard, 
after which the error flags were cleared and the motors operated 
normally. Both motors were retained for use in the final assembly.

\hmsubsubsection{Thermal Protection at Scan Pose}{UMA}{}
A practical issue identified during integration testing concerned 
the elbow motor (M3, XM430-W350). The system spends a significant 
proportion of its operating time holding the arm at a fixed 
\texttt{SCAN\_POSE} configuration, in which the camera is positioned 
above the workspace for ball detection (see Section~8.7). 
At this pose, the elbow motor must continuously resist the 
gravitational torque produced by the forearm, the wrist, the 
gripper, and the camera. Holding torque draws current 
continuously, and the XM430-W350 was observed to draw 
approximately 0.47~A while stationary at this pose. Over extended 
periods, this current causes the motor to heat up, which can 
trigger an overheating error flag and cause the motor to shut 
down mid-operation.

To address this, a reduced current limit of 300~mA was
implemented and tested on the elbow motor whenever the arm
was parked at \texttt{SCAN\_POSE}. The intention was to cap
the steady-state current and thereby reduce heat generation
while still providing sufficient torque to hold the arm in
position under nominal load.

\textbf{Outcome.}
Testing showed that the reduced current limit did not
produce a measurable improvement in motor temperature
during extended operation. The thermal load on the elbow
motor is fundamentally determined by the gravitational
torque from the weight of the forearm, camera, and gripper
assembly, which the motor must resist regardless of the
current limit setting. A current limit can only prevent
the motor from drawing excessive current during transient
loads; the steady-state holding current is already below
the stall current and cannot be reduced further without
allowing the arm to droop.

\textbf{Conclusion.}
The thermal protection logic was therefore removed from the
deployed system. The most effective mitigation identified
would be to reduce the mass of the L\textsubscript{3}
assembly (the wrist-to-claw link, camera, and gripper),
which would directly decrease the gravitational torque on
the elbow and the corresponding holding current. During
the project's testing sessions the motor has not exhibited
thermal shutdown; however, extended continuous operation
could still pose a risk and should be monitored if the
system is operated for longer periods.